% !TeX program = xelatex
% !TeX encoding = UTF-8
\documentclass{MathorCupmodeling}

% === 额外宏包 ===
\usepackage{algorithm2e}
\usepackage{float}
\usepackage[section]{placeins}
\usepackage{needspace}

% === TikZ ===
\usepackage{tikz}
\usetikzlibrary{arrows.meta,positioning,shapes.geometric,fit,calc}
\usepackage{longtable}

% === 竞赛信息 ===
\bianhao{[队伍编号]}
\tihao{[题号]}
\timu{[论文标题]}
\keyword{[关键词1]；[关键词2]；[关键词3]}

% === 引用上标格式 ===
\makeatletter
\def\@cite#1#2{\textsuperscript{[{#1\if@tempswa , #2\fi}]}}
\makeatother

\begin{document}

% === 摘要 ===
\begin{abstract}

[中文摘要内容：问题概述 + 每个子问题的方法和数值结果 + 结论，400-600字]

\end{abstract}

% === 目录 ===
\newpage
\tableofcontents
\newpage

% === 正文 ===
\input{sections/1_restatement}
\input{sections/2_analysis}
\input{sections/3_assumptions}
\input{sections/4_symbols}
\input{sections/5_problem1}
\input{sections/6_problem2}
\input{sections/7_problem3}
\input{sections/8_evaluation}

% === 参考文献 ===
\phantomsection
\addcontentsline{toc}{section}{参考文献}
\begin{thebibliography}{99}
% \bibitem{ref1} 作者. 标题[J]. 期刊, 年份.
\end{thebibliography}

% === 附录 ===
\appendix
\ctexset{section={format={\zihao{-4}\heiti\raggedright}}}
\begin{center}
  \heiti\zihao{4} 附\hspace{1pc}录
\end{center}
\input{sections/A_code}

\end{document}
